\appendix
\section{Resolution robustness and sign cancellation}

This appendix explores an alternative signed-aggregate concentration definition
to illustrate how different choices in aggregation order affect resolution
sensitivity. All results and inferences in the main text use the abs-before-sum
definition specified in \S2.5; this appendix does not alter those conclusions.

\paragraph{Alternative signed-aggregate construction.}
To illustrate resolution behavior under different aggregation choices, we
consider a signed-aggregate variant that applies the absolute value after
summing child-pixel contributions within parent pixels. We recomputed this
alternative Top-20\% statistic while varying the child HEALPix resolution over
$\mathrm{NSIDE}_{\mathrm{child}} \in \{64,128,256\}$,
holding the parent grid fixed at
$\mathrm{NSIDE}_{\mathrm{parent}} = 64$.
Across this range, the signed-aggregate concentration varies modestly, from
$C_{0.20}^{\mathrm{signed}} = 0.572$ at
$\mathrm{NSIDE}_{\mathrm{child}} = 64$
to $0.540$ at $256$.
Because this construction operates on a fixed parent grid with varying child
resolution, its numerical values are not expected to match the single-resolution
$C_{0.20}$ values reported in the main text.

\paragraph{Sign cancellation under refinement.}
To quantify why different aggregation orders exhibit different resolution
dependence, we compute a sign-cancellation diagnostic
\[
\mathrm{cancel\_ratio}
= \frac{\sum_P |S_P|}{\sum_c |s_c|},
\]
where $s_c = N_c \tilde{\varpi}_c$ denotes signed child-pixel contributions and
$S_P = \sum_{c \in P} s_c$ their aggregation into parent pixels.
A value of $\mathrm{cancel\_ratio} < 1$ indicates cancellation of opposite-sign
contributions within parent pixels.

We find $\mathrm{cancel\_ratio} = 1.00$ at
$\mathrm{NSIDE}_{\mathrm{child}} = 64$, decreasing to $0.64$ at $128$ and
$0.40$ at $256$.
This behavior indicates increasing within-parent sign cancellation at finer
resolution. Abs-before-sum statistics applied at a single resolution avoid
this effect by construction, which motivates the choice in the main text.
The signed-aggregate alternative offers complementary insight into the
spatial structure of the residual field but is not used for inference in
this work.

\begin{figure}[t]
\centering
\includegraphics[width=0.75\textwidth]{figures/lorenz_curve_signed.pdf}
\caption{\textbf{Alternative signed-aggregate Lorenz curve.}
Lorenz-style curve using the signed-aggregate construction with
$\mathrm{NSIDE}_{\mathrm{child}} = 128$ aggregated to
$\mathrm{NSIDE}_{\mathrm{parent}} = 64$, yielding $C_{0.20} = 0.56$.
This construction illustrates the effect of sign cancellation discussed
in this appendix; all inferences in the main text use the single-resolution
abs-before-sum definition (Figure~2).}
\label{fig:lorenz_signed_appendix}
\end{figure}
